\documentclass[11pt]{article}
\usepackage{fontspec}
\setmainfont{Times New Roman}
\setsansfont{Arial}
\setmonofont{Consolas}
\usepackage{amsmath,amssymb,mathtools}
\usepackage{geometry}
\usepackage{microtype}
\geometry{a4paper,margin=2.28cm}
\setlength{\parindent}{1.15em}
\setlength{\parskip}{0pt}
\makeatletter
\renewcommand\footnotesize{\@setfontsize\footnotesize{7.5}{8.5}\selectfont}
\makeatother
\providecommand{\mG}{\mathfrak{G}}
\providecommand{\mA}{\mathfrak{A}}
\providecommand{\mE}{\mathfrak{E}}
\providecommand{\mH}{\mathfrak{H}}
\providecommand{\mR}{\mathfrak{R}}
\providecommand{\mo}{\mathfrak{o}}
\providecommand{\iso}{\simeq}
\emergencystretch=120em
\allowdisplaybreaks
\sloppy
\hbadness=10000
\vbadness=10000
\hfuzz=2pt
\vfuzz=10pt
\raggedbottom
\begin{document}
% Unit: Paper34 section06 v001. Direct Interslavic Latin translation from audited German source lines 334--404, segments P34-S0074--P34-S0084. Source scan pages 13--15 retained as witnesses, and Zenodo record 20822156 was rechecked for this unit.
\setcounter{footnote}{11}
\section*{34. Hiperkompleksne veličiny i teorija predstavjenj}
\emph{Prodolženje: glava I, §§5--7, i glava II, §§8--14.}

\subsection*{§ 6. Moduly vzhodno k tělu. Hiperkompleksne sistemy}

Skoro trivialny primjer popolno reducibilnyh grup s operatorami daju moduly s konečnoju bazoju vzhodno k (ne nužno komutativnomu) tělu.

Nehaj $\mG$ bude pravy $K$-modul, i nehaj jediničny element $e$ těla $K$ bude Jednočasno jediničny operator: $ae=a$ za $a$ iz $\mG$. Vsaky podmodul $aK$, generovany elementom $a$, jest prosty, imenno jednak modulu generovanomu libovoljnym jego nenuljevym elementom. Zato, ako $\mH$ jest podmodul generovany někojimi elementami, a $a$ jest element ne naležeči k $\mH$, togda $\mH\cap aK=\mE$, i zato $\mH+aK$ jest direktna suma. Tako, izhodeči iz libovoljnogo bazovogo elementa $a_1$, možno vsečas dodavati nove bazove elementy $a_2,a_3,\ldots$ i dobiti $\mG$ kako direktnu sumu
\[
        \mG=a_1K+a_2K+\cdots+a_nK.
\]
\emph{Zato $\mG$ jest popolno reducibilna.}

Čislo $n$, to jest dolgost kompozicijnogo rjada, nazyvaje se \emph{rangom} $\mG$ vzhodno k $K$. Baza $a_1,\ldots,a_n$ jest linearno nezavisna zaradi jednoznačnogo predstavjenja elementov $\mG$ (i zato, že $e$ bylo predpoloženo kako jediničny operator).

Vsaky podmodul $\mH$ jest direktny sumand:
\[
        \mG=\mH+\mR=(h_1K+\cdots+h_sK)+(r_1K+\cdots+r_{n-s}K),
\]
ili: \emph{vsaku linearno nezavisnu bazu $\mH$ možno dopolniti do linearno nezavisnoj bazy $\mG$.} Elementy $r_i$ možno čak izbrati iz izhodnyh bazovyh elementov $a_i$.

Nehaj $(a_1,\ldots,a_n)$ i $(c_1,\ldots,c_n)$ budut linearno nezavisne bazy. Togda
\[
        c_k=\sum_i a_i\pi_{ik},
\]
ili, zapisano v matricah,
\[
        (c_1,\ldots,c_n)=(a_1,\ldots,a_n)P,
        \qquad
        P=\begin{pmatrix}
        \pi_{11}&\cdots&\pi_{1n}\\
        \vdots&&\vdots\\
        \pi_{n1}&\cdots&\pi_{nn}
        \end{pmatrix}.
\]
Obratno,
\[
        (a_1,\ldots,a_n)=(c_1,\ldots,c_n)Q.
\]
Iz togo slěduje
\[
        (a_1,\ldots,a_n)=(a_1,\ldots,a_n)PQ,
        \qquad
        (c_1,\ldots,c_n)=(c_1,\ldots,c_n)QP,
\]
a zato, poneže $a_i$ i $c_i$ sut linearno nezavisne,
\[
        PQ=QP=E.
\]

Specialne $K$-moduly, ktore sut Jednočasno kolca, sut ``hiperkompleksne sistemy''.

Kolco $\mo$, ktore jest Jednočasno pravy modul vzhodno k komutativnomu polju $K$, nazyvaje se \emph{hiperkompleksnym sistemom vzhodno k $K$} (v literaturi takože ``algebra nad $K$''), ako:
\begin{enumerate}
\item rang jest konečny (nehaj linearno nezavisna baza bude $u_1,\ldots,u_n$);
\item $ab\cdot x=a\cdot bx=ax\cdot b$. To izraža se tako: $K$ jest komutativno svęzano s $\mo$;
\item jediničny element $e$ polja $K$ jest Jednočasno jediničny operator: $ae=a$ za $a$ iz $\mo$.
\end{enumerate}
Zaradi
\[
        \Bigl(\sum_i u_i\alpha_i\Bigr)\Bigl(\sum_k u_k\beta_k\Bigr)
        =\sum_{i,k}(u_i u_k)(\alpha_i\beta_k),
\]
sistem jest jednoznačno opreděljen, kogda, kromě $K$, dana jest \emph{tablica množenja}, ktora ukazuje, kako vsaky produkt $u_i u_k$ izražaje se črěz $u_j$:
\[
        u_i u_k=\sum_j u_j\gamma^j_{ik}.
\]
Elementy $\gamma^j_{ik}$ muset zadovoliti znane relacije, ktore slědujut iz asociativnogo zakona.

Ako $\mo$, kako často budemo predpolagati, imaje jediničny element $e$ ($ea=ae=a$ za vse $a$), togda elementy $x$ polja $K$ možno identificirati s $ex$, to jest smatrjati $K$ kako podpolje v $\mo$. Hiperkompleksny sistem možno togda takože opisati kako \emph{kolco konečnogo ranga vzhodno k polju ležečemu v centru.}

Primjer hiperkompleksnogo sistema jest \emph{grupno kolco} konečnoj grupy: za bazove elementy $u_i$ berut se elementy konečnoj grupy, a za množenje -- grupno množenje. Polje $K$ pri tom jest libovoljno.

\footnotetext[12]{Uslovje 2 govori, že množenje elementom iz $K$ davaje homomorfizm za $\mo$, kogda $\mo$ razsmatra se i kako lěvy $\mo$-modul, i kako pravy $\mo$-modul. Posrědstvom 3 ty homomorfizmy sut čak izomorfizmy.}
\end{document}
